\section{Metodología}

Para el análisis de cada uno de los GRBs estudiados en este trabajo (GRB 080916C y GRB 251118A) se siguió el flujo de trabajo estándar recomendado por el \textit{Fermi Science Support Center} para el análisis de datos del LAT \cite{FSSCCicerone}, el cual combina un análisis \textit{binned} (para caracterizar la fase prompt, de corta duración) y un análisis \textit{unbinned} de máxima verosimilitud (\textit{unbinned likelihood analysis}) para la emisión extendida de alta energía. Todas las gráficas de resultados de este informe se elaboraron con Python (\texttt{astropy}, \texttt{numpy} y \texttt{matplotlib}), excepto las de los ajustes espectrales en XSPEC, que se obtuvieron con la salida gráfica del propio programa. A continuación se presenta con más detalle el paso a paso que se realizó para el análisis de estos dos GRBs.

\subsection{Obtención de los datos}

El primer paso consiste en obtener los datos crudos del LAT a través del \textit{LAT Photon, Event, and Spacecraft Data Query} del FSSC. Para esto, es necesario que especifiquemos los siguientes parámetros de búsqueda: la posición del GRB en coordenadas ecuatoriales (RA, Dec), el radio de la región de interés (\textit{Region of Interest}, ROI), el tiempo de inicio y finalización del intervalo de interés en Tiempo del Instrumento (\textit{Mission Elapsed Time}, MET), y el rango de energía mínima y máxima en MeV. Como resultado de esta consulta se obtienen dos archivos en formato FITS: el archivo de eventos (\textit{Event file}, FT1), que contiene la lista de los eventos reconstruidos como candidatos a fotón —``evento'' en el sentido de la sección 2.2— con su tiempo de llegada, energía y dirección, y el archivo de la nave espacial (\textit{Spacecraft file}, FT2), que registra hacia dónde apuntaba el telescopio y con qué orientación durante todo el intervalo. Esa información de apuntamiento, combinada con las IRFs, es la que después permite calcular la exposición.

\subsection{Selección inicial de eventos (\texttt{gtselect})}

Aunque los archivos ya vienen recortados por la consulta, esa descarga es solo un recorte grueso: un cono espacial y una ventana de tiempo y energía sobre los datos extendidos, que traen todas las clases de eventos mezcladas. Los cortes finos se hacen en esta etapa con \texttt{gtselect}, que aplica criterios espaciales, temporales, energéticos y de calidad sobre el archivo original. Aquí se utiliza un radio de ROI relativamente amplio (típicamente de 15$^{\circ}$, holgado también frente a la incertidumbre de la posición reportada en la circular) y no se fija aún la posición exacta del centro, con el fin de contar con suficientes fotones para una localización precisa posterior de la fuente. Se aplica además un corte en el ángulo cenital (\texttt{zmax}, que en este trabajo fue de 100$^{\circ}$) para excluir fotones provenientes del limbo terrestre, y es en este paso —no en el servidor— donde se elige la clase de eventos (\texttt{evclass}) según la fase de análisis: la clase \texttt{TRANSIENT} para la emisión prompt, de mayor sensibilidad a GRBs cortos y débiles pero mayor contaminación de rayos cósmicos o fotones generados por interacción con la nave, o la clase \texttt{SOURCE}, más estricta en la discriminación del fondo, pero más sensible para GRBs de cientos de segundos.

Las llamadas completas a cada herramienta, con todos sus parámetros, se recogen en el Anexo~\ref{anexo:comandos} para no interrumpir la lectura; en el texto se citan solo los parámetros que definen cada decisión. La de \texttt{gtselect} es la del Anexo~\ref{an:gtselect}.

\subsection{Curva de luz preliminar (\texttt{gtbin})}

Con el archivo filtrado se genera una curva de luz preliminar utilizando \texttt{gtbin} en modo \texttt{LC} (\textit{Light Curve}), agrupando los eventos en bins temporales de duración fija (\texttt{dtime}). La curva se grafica con Python, con el tiempo en el eje horizontal y el número de conteos por bin en el eje vertical, lo que permite identificar visualmente la duración aproximada del estallido y su estructura temporal (Anexo~\ref{an:lc}).

\subsection{Localización de la fuente (\texttt{gtfindsrc})}

Con el objetivo de refinar la posición del GRB proporcionada inicialmente (usualmente proveniente de la localización del GBM o de una circular GCN), se utiliza \texttt{gtfindsrc}. Esta herramienta realiza un ajuste iterativo que maximiza la verosimilitud de un modelo de fuente puntual única en función de su posición, partiendo de una posición semilla y usando las IRFs (sección 2.2) correspondientes a la fase de emisión analizada. El resultado es una posición mejorada (RA, Dec) que se utiliza en los pasos subsiguientes del análisis (Anexo~\ref{an:gtfindsrc}).

\subsection{Mapa de conteos (\texttt{gtbin})}

Con la posición refinada se genera un mapa de conteos bidimensional (\texttt{CMAP}) mediante \texttt{gtbin}, en coordenadas celestes ecuatoriales (J2000). El mapa se grafica con Python para inspeccionar la distribución espacial de los eventos alrededor de la posición del GRB y confirmar visualmente que hay un exceso de conteos en la posición de la fuente.

\subsection{Definición del modelo de la región de interés}

Antes de proceder con el ajuste espectral, es necesario definir un modelo de la región de interés (ROI) en formato XML. Este modelo incluye, como mínimo, tres componentes: (i) una fuente puntual (\textit{Point Source}) ubicada en la posición refinada del GRB, con un espectro tipo ley de potencia (\texttt{PowerLaw2}); (ii) una componente de emisión difusa galáctica, que representa la emisión de rayos gamma producida por la interacción de rayos cósmicos con el gas interestelar de la Vía Láctea, modelada a partir del mapa estándar \texttt{gll\_iem\_v07.fits}, construido con el código de propagación de rayos cósmicos GALPROP (de ahí el nombre con el que esta componente aparece en los resultados); y (iii) una componente isotrópica, que representa el fondo difuso extragaláctico y el fondo residual de partículas cargadas mal clasificadas. Los archivos de fondo difuso estándar se obtienen a través del \textit{Model Editor} del FSSC, herramienta que genera automáticamente el archivo XML con la estructura y los parámetros iniciales adecuados para cada componente.

Fermitools permite escribir la ley de potencia de la fuente puntual de dos maneras, que describen la misma forma espectral y solo se diferencian en cómo se fija su normalización \cite{FSSCCicerone}. En \texttt{PowerLaw}, la normalización es un prefactor, es decir, el número de fotones por unidad de energía, de área y de tiempo a una energía de referencia (en este trabajo, 100~MeV). En \texttt{PowerLaw2}, en cambio, la normalización es directamente el flujo de fotones integrado entre dos energías (aquí, de 100~MeV a 300~GeV), de modo que el ajuste entrega el flujo y su error sin cálculos adicionales. Por esa razón el análisis se hizo con \texttt{PowerLaw2}, y \texttt{PowerLaw} se usó solo como control.

\subsection{Cálculo de la contribución difusa (\texttt{gtdiffrsp})}

Una vez definido el modelo XML de la región de interés, se ejecuta \texttt{gtdiffrsp}, la cual calcula la respuesta del instrumento a las componentes difusas del modelo para cada fotón individual: la integral del modelo difuso, convuelta con la PSF y el área efectiva, evaluada en la dirección, la energía y el tiempo de llegada de ese fotón (por eso el cálculo es fotón por fotón). El resultado se agrega como nuevas columnas (\texttt{DIFRSP0}, \texttt{DIFRSP1}, \dots) en el archivo de eventos. Este paso es indispensable para que el ajuste de máxima verosimilitud posterior pueda distinguir correctamente entre la contribución de la fuente puntual y la de los fondos difusos (Anexo~\ref{an:gtdiffrsp}).

\subsection{Livetime cube y mapa de exposición}

Posteriormente se calcula el \textit{livetime cube} con \texttt{gtltcube}, el cual almacena el tiempo efectivo de observación acumulado para cada posición del cielo, en función del ángulo con que esa posición incide respecto al eje del instrumento, a lo largo de todo el intervalo analizado (es una propiedad de cómo apuntó el telescopio, no de la fuente). A partir de este cubo se genera, mediante \texttt{gtexpmap}, un mapa de exposición: la integral del área efectiva por el tiempo efectivo de observación en cada punto y energía de la ROI, con unidades de cm$^2\,$s. La exposición es el factor que convierte flujos en conteos esperados, y por eso es indispensable para el ajuste espectral.

\subsection{Análisis binned: espectro y ajuste en XSPEC}

Para la fase prompt se realiza un análisis \textit{binned} tradicional: al ser corta y brillante, hay conteos suficientes para agruparlos por energía y ajustar el espectro acumulado (la emisión extendida, con pocos fotones repartidos en mucho tiempo, se trata con el método \textit{unbinned} de la sección 3.10). En primer lugar se genera el espectro de conteos observado con \texttt{gtbin} en modo \texttt{PHA1}, que produce un único espectro acumulado sobre toda la ventana en formato PHA (el formato estándar de espectros en astronomía de rayos X), definiendo los bordes de energía en escala logarítmica. A continuación se genera la matriz de respuesta espectral (RSP) correspondiente con \texttt{gtrspgen}, utilizando el método de cálculo \texttt{PS} (\textit{Point Source}), recomendado sobre el método \texttt{GRB} por ofrecer una mayor precisión, en particular para estallidos de mayor duración; para estallidos cortos la diferencia entre ambos métodos, tanto en resultados como en tiempo de ejecución, es despreciable. La matriz conecta dos mallas de energía distintas: la de la energía verdadera del fotón incidente, que se define en \texttt{gtrspgen}, y la de los canales del archivo PHA, que corresponden a la energía reconstruida. Como la resolución en energía del instrumento es finita, conviene que la primera malla sea más fina que la segunda (en este trabajo, 100 bins frente a 30 canales): así el número esperado de fotones en cada canal se calcula con mayor precisión (Anexo~\ref{an:binned}).

Con el espectro y su respuesta ya generados, se procede al ajuste en \texttt{XSPEC}, empleando un modelo de ley de potencia simple (\texttt{pow}). Al definir este modelo, \texttt{XSPEC} solicita como valores iniciales el índice fotónico (\texttt{PhoIndex}) y la normalización (\texttt{norm}) de la ley de potencia. El ajuste se realiza minimizando la estadística de Cash (\texttt{cstat}), apropiada para el régimen de bajo número de conteos característico de la emisión de alta energía del LAT, y se inspeccionan visualmente tanto el espectro ajustado como los residuos normalizados por su incertidumbre (Anexo~\ref{an:xspec}).

\subsection{Análisis unbinned: ajuste de máxima verosimilitud (\texttt{gtlike})}

Finalmente se analiza la emisión extendida; extendida en el tiempo: en muchos GRBs los fotones de alta energía no llegan en el pico del estallido, sino repartidos a lo largo de cientos de segundos. Para esta fase se usa el análisis de máxima verosimilitud \textit{unbinned}, en el que la verosimilitud se evalúa fotón por fotón en lugar de agrupar los eventos en bins; al no agrupar no se pierde información, y por eso es el método apropiado cuando los fotones detectados son pocos. El intervalo temporal es más amplio que el de la fase prompt y se emplea la clase \texttt{SOURCE} con su IRF (\texttt{P8R3\_SOURCE\_V3}), más estricta en la discriminación del fondo. El flujo de trabajo repite, sobre este nuevo intervalo, la selección de eventos (\texttt{gtselect}), el filtrado de calidad mediante \texttt{gtmktime} —aplicando la máscara \texttt{(DATA\_QUAL>0)\&\&(LAT\_CONFIG==1)}, que es la que el FSSC recomienda en el Cicerone para quedarse solo con datos de buena calidad y el instrumento en configuración científica \cite{FSSCCicerone}—, el cálculo de la respuesta difusa (\texttt{gtdiffrsp}, ahora con la IRF de \texttt{SOURCE}), y la generación del livetime cube y del mapa de exposición correspondientes a este nuevo intervalo, también con esa IRF (Anexo~\ref{an:extendida}).

Con todos estos archivos generados (evento filtrado, livetime cube, mapa de exposición y modelo XML), se ejecuta el ajuste de máxima verosimilitud propiamente dicho mediante la herramienta \texttt{gtlike}, en modo \texttt{UNBINNED} (el método fotón por fotón recién descrito) y con el optimizador \texttt{MINUIT} \cite{James1975}. Esta herramienta ajusta simultáneamente los parámetros libres de todas las componentes del modelo XML (la fuente puntual del GRB, la componente galáctica y la componente extragaláctica) maximizando la función de verosimilitud. La verosimilitud es la probabilidad de haber observado exactamente los datos que se observaron si el modelo fuera cierto, vista como función de los parámetros del modelo: ajustar por máxima verosimilitud es buscar los valores de los parámetros que hacen esa probabilidad lo más alta posible. En el caso \textit{unbinned} se construye como el producto, sobre todos los fotones observados, de la densidad de conteos que el modelo predice en la energía, posición y tiempo de llegada de cada uno, multiplicado por el término de Poisson del número total de conteos esperado \cite{Mattox1996} (la llamada, en el Anexo~\ref{an:gtlike}).

Como resultado de este ajuste, \texttt{gtlike} produce un archivo de salida (\texttt{.dat}) que reporta, para cada componente del modelo, los valores ajustados de sus parámetros espectrales junto con su incertidumbre estadística, el número de fotones predichos (\texttt{Npred}), el flujo integrado en el rango de energía analizado, y en el caso de la fuente puntual, el valor del estadístico de prueba (\textit{Test Statistic}, TS), definido como $\mathrm{TS} = -2\,\ln(\mathcal{L}_{0}/\mathcal{L}_{1})$, donde $\mathcal{L}_{0}$ y $\mathcal{L}_{1}$ son las verosimilitudes máximas del modelo nulo (sin la fuente) y del modelo que incluye la fuente puntual, respectivamente \cite{Mattox1996}. Este valor de TS es el criterio estándar para cuantificar la detección de un GRB con el LAT. El TS no es en sí la significancia: por el teorema de Wilks, cuando el modelo agrega un parámetro libre, la significancia aproximada en unidades de $\sigma$ es $\sqrt{\mathrm{TS}}$, y de ahí viene la regla práctica de que $\mathrm{TS} \gtrsim 25$ corresponde aproximadamente a una detección de $5\sigma$. Con pocos fotones esta aproximación es solo orientativa.